\documentclass[12pt, oneside, final, letterpaper, stu]{huthesis}


%%%% Paquetes %%%%%%%%%%%%%%%%%
\usepackage{iftex}
\ifLuaTeX
  \usepackage{fontspec} % LuaLaTeX: Latin Modern en Unicode (acentos y signos ¿ ¡ correctos)
\else
  \usepackage[utf8]{inputenc}
  \usepackage[T1]{fontenc}
  \usepackage{lmodern}
\fi
\usepackage{verbatim}
%\usepackage[latin1]{inputenc}
%\inputencoding{latin1}
%\usepackage[spanish]{babel}
%\usepackage[latin1]{inputenc}

%\usepackage[latin1]{inputenc}
\usepackage{graphicx}
\usepackage{xcolor,colortbl}
% Compatibilidad: colortbl (2026) usa \insert@pcolumn, que solo existe en array >= 2024.
% Si la versión instalada de array es anterior, se define como \insert@column.
\makeatletter
\providecommand*{\insert@pcolumn}{\insert@column}
\makeatother
\usepackage{rotating}
\usepackage{pdflscape}
\usepackage{graphicx}
\usepackage{amsmath}
\usepackage{amssymb}
\usepackage{amssymb}
\usepackage{amsmath}
\usepackage{mathrsfs}
\usepackage[spanish]{babel}
%\usepackage[print]{pdfscreen}
% \usepackage{ae} % ELIMINADO: forzaba fuentes OT1 y anulaba fontenc T1 (acentos no copiables en el PDF)
\usepackage{array}
\usepackage{textcomp}
\usepackage{amsmath}
\usepackage{amssymb}
\usepackage{subfig}
\usepackage{mathrsfs}
%\usepackage{matlab}
\usepackage{hyperref}
\usepackage{Portada}
%CORREGIR LA METODOLOGIA
\usepackage{multirow}
\usepackage{fancyhdr}
%\usepackage{algoritmo}
\usepackage{epsfig}
\usepackage{float}
\usepackage{csquotes}
\usepackage{colortbl}
\usepackage{longtable}
\usepackage{pgfgantt}
\usepackage{booktabs}
\usepackage{tabularx}
\usepackage{makecell}
\usepackage{array}
\usepackage{ragged2e}
\usepackage{pdflscape}
\usepackage{microtype}
\usepackage{rotating}
\usepackage{adjustbox}
\usepackage{amssymb}
\usepackage{pifont}
\usepackage{caption}
% --- Resaltado de los cambios de la versión 2 ---
% \nuevo{...} resalta en amarillo el texto nuevo o corregido. Requiere LuaLaTeX
% (configurado en .latexmkrc). Para la versión limpia, cambiar \resaltartrue por \resaltarfalse.
\newif\ifresaltar
\resaltartrue
\ifLuaTeX
  \usepackage{luacolor}
  \usepackage{lua-ul}
  \LuaULSetHighLightColor{yellow}
\fi
\makeatletter
\ifresaltar
  \ifLuaTeX
    \DeclareRobustCommand{\nuevo}[1]{\highLight{#1}}
    \newcommand{\nuevobib}{\@highLight}   % interruptor para las entradas nuevas de la bibliografía
  \else
    % Con pdfLaTeX no se puede resaltar texto con citas y fórmulas; se usa color como respaldo.
    \DeclareRobustCommand{\nuevo}[1]{\textcolor{orange!80!black}{#1}}
    \newcommand{\nuevobib}{\color{orange!80!black}}
  \fi
\else
  \DeclareRobustCommand{\nuevo}[1]{#1}
  \newcommand{\nuevobib}{}
\fi
\makeatother
\usepackage[backend=biber, style=apa]{biblatex}
\addbibresource{bibliografia.bib}
% Las referencias agregadas o corregidas en la versión 2 llevan keywords = {nuevo} en el .bib
\AtEveryBibitem{\ifkeyword{nuevo}{\nuevobib}{}}

% APA 7 en español: "y" en lugar de "&" entre los últimos autores
\DeclareDelimFormat[bib,biblist]{finalnamedelim}{\addspace y\space}
\DeclareDelimFormat[textcite]{finalnamedelim}{\addspace y\space}
\DeclareDelimFormat[parencite]{finalnamedelim}{\addspace y\space}
% Sin coma antes de "y" en la lista de referencias
\makeatletter
\renewcommand*{\apablx@ifnotfinalname}{%
  \ifboolexpr{test {\ifnumless{\value{listcount}}{\value{liststop}-1}} or test \ifmorenames}}
\makeatother



\newcommand{\YES}{\ding{51}}   % relación directa
\newcommand{\PAR}{$\triangle$} % relación parcial
\newcommand{\NON}{---}         % sin relación


%%%%%%%%%%%
\usepackage{color}
\definecolor{gray97}{gray}{.97}
\definecolor{gray75}{gray}{.75}
\definecolor{gray45}{gray}{.45}

\usepackage{listings}
\lstset{ frame=Ltb,
framerule=0pt,
aboveskip=0.5cm,
framextopmargin=3pt,
framexbottommargin=3pt,
framexleftmargin=0.4cm,
framesep=0pt,
rulesep=.4pt,
backgroundcolor=\color{gray97},
rulesepcolor=\color{black},
%
stringstyle=\ttfamily,
showstringspaces = false,
basicstyle=\small\ttfamily,
commentstyle=\color{gray45},
keywordstyle=\bfseries,
% 
numbers=left,
numbersep=15pt,
numberstyle=\tiny,
numberfirstline = false,
breaklines=true,
}


\lstdefinestyle{consola}
{basicstyle=\scriptsize\bf\ttfamily,
backgroundcolor=\color{gray75},
}

\lstdefinestyle{C}
{language=C,
}
%%%%%%%%%%%


\usepackage{listings}
\usepackage{color}
\usepackage{multicol}

%\usepackage[latin1]{inputenc}
%\usepackage[utf8]{inputenc}

\hypersetup{backref=false, colorlinks=false}
\spanishdecimal{.}
\DeclareGraphicsExtensions{.pdf}
\def\tablesp{\def\baselinestretch{1.1}\large\normalsize}
\def\subtema{\subsubsection*}

\usepackage{enumitem,amssymb}
\newlist{todolist}{itemize}{2}
\setlist[todolist]{label=$\square$}

\usepackage{xcolor,colortbl}
\definecolor{verde}{rgb}{0,0.9,0}

% Numeración de tablas sin número de sección
\counterwithout{table}{section}

%----------------------------------
\hyphenation{  }
     % ---- Se••••••••paraciones de palabras
% --- Conceptos previos --- Apéndice A
\newcommand{\Fc}{$\textbf{F}$}

%%Definiciones
\newcommand{\vxd}{\textbf{x}}
\newcommand{\vyd}{\textbf{y}}
\newcommand{\vzd}{\textbf{z}}
\newcommand{\vsd}{\textbf{s}}
\newcommand{\ved}{\textbf{e}}

%ecuaciones
\newcommand{\vx}{\mathbf{x}}
\newcommand{\vy}{\mathbf{y}}
\newcommand{\vz}{\mathbf{z}}
\newcommand{\vs}{\mathbf{s}}
\newcommand{\ve}{\mathbf{e}}
     % ---- Definiciones matemticas


%------------ Informacin de la Tesis
\title{\Large \nuevo{Efecto de una plataforma web en el tiempo de búsqueda y la confianza al rentar vivienda: estudiantes de la UNSIS}}	%---> Titulo
\author{\nuevo{Mayra Castellanos Pacheco} \\ \nuevo{Elisa Sánchez Martínez}}												%---> Autor
\usepackage[spanish]{babel}


\degreemonth{Noviembre}																		%---> Mes
\degreeyear{2026}																		%---> Ao
\degree{Licenciado en Informática}															%---> Grado
\institute{Instituto de Informática}														%---> Nombre del Instituto
\department{Instituto de Informática}														%---> Nombre del departamento
\university{Universidad de la Sierra Sur}													%---> Nombre universidad
\place{Miahuatlán de Porfirio Díaz Oaxaca, México.}											%---> Lugar
% TODO: confirmar que el reglamento de titulación permite cuatro directores/asesores.
\advisor{Asesores}{M.C.C. Eliezer Alcázar Silva \\ M.C.C. Lirio Ruiz Guerra \\ M.F. Máximo Jorge Saavedra García \\ M.T.E. Everardo de Jesús Pacheco Antonio }	
%Algoritmo
\definecolor{dkgreen}{rgb}{0,0.6,0}
\definecolor{gray}{rgb}{0.5,0.5,0.5}
\definecolor{mauve}{rgb}{0.58,0,0.82}
\newlength\mylength
\newcolumntype{C}[1]{>{\centering\arraybackslash}p{#1}}
\setlength\mylength{\dimexpr\textwidth-5\arrayrulewidth-8\tabcolsep}

\lstset{
	frame=tb,
	language=C++,
	aboveskip=3mm,
	belowskip=3mm,
	numbers=left,
	showstringspaces=false,
	columns=flexible,
	basicstyle={\small\ttfamily},
	numberstyle=\tiny\color{gray},
	keywordstyle=\color{blue},
	commentstyle=\color{dkgreen},
	stringstyle=\color{mauve},
	breaklines=true,
	frame=leftline,
	breakatwhitespace=true,
	tabsize=3
}

%%%%%%%%%%%%%%%%%%%%%%%%%%%%%%%%%%%%%%%%%%%%%%%%%%%%%%%%%%%%%%%%%%%%%%%%%%%%%555

%

% --- Configuración para numeración de secciones ---
\renewcommand{\thesection}{\thechapter.\arabic{section}}
\renewcommand{\thesubsection}{\thesection.\arabic{subsection}}
\renewcommand{\thesubsubsection}{\thesubsection.\arabic{subsubsection}}

\setlength{\headheight}{14.5pt} % evita la advertencia de fancyhdr

\begin{document}


    % Configuración de encabezados
    \pagestyle{fancy}
    \fancyhf{}
    \fancyhead[L]{\slshape \leftmark}
    \fancyfoot[C]{\thepage}
    \renewcommand{\headrulewidth}{0.4pt}
    \renewcommand{\chaptermark}[1]{\markboth{Capítulo \thechapter: #1}{}}
    \renewcommand{\sectionmark}[1]{}

    \renewcommand{\tablename}{Tabla}
    \renewcommand{\listtablename}{Índice de tablas}
    \maketitle

    %--- Páginas preliminares (comentadas mientras estén vacías) ---
%    \begin{dedication}
%        \vspace*{\fill}
%        \begin{flushright}
%            \textit{Dedicatoria aquí...}
%        \end{flushright}
%        \vspace*{\fill}
%    \end{dedication}
%
%    \begin{acknowledgments}
%        \begin{itemize}
%            \item A alguien
%            \hspace*{\fill} \item A otro alguien
%            \item A otro alguien
%        \end{itemize}
%    \end{acknowledgments}
%
%    \begin{resumen}
%        Resumen en español aquí...
%    \end{resumen}
%
%    \begin{abstract}
%        Abstract in English here...
%    \end{abstract}

%--- Índices y páginas preliminares ---
    \pagenumbering{roman} 
    \setcounter{page}{0}

    \tableofcontents
    \cleardoublepage

%    \listoffigures      % sin figuras por ahora
%    \cleardoublepage

%    \listoftables       % sin tablas por ahora
%    \cleardoublepage

 %  % --- Listado de Abreviaturas y Acrónimos ---
    \chapter*{Listado de Abreviaturas y Acrónimos}
    \addcontentsline{toc}{chapter}{Listado de Abreviaturas y Acrónimos} 
    
    \begin{description}
        \item[IA] Inteligencia Artificial
        \item[DCL] Deterioro Cognitivo Leve
        \item[API] Interfaz de Programación de Aplicaciones
    \end{description}
%    \cleardoublepage

 %   \section
{Glosario}
\label{cap:glosario}

El presente glosario reúne los términos técnicos y conceptuales de mayor 
recurrencia en esta investigación, organizados alfabéticamente. Cada 
definición incluye la fuente de referencia que le da sustento dentro del 
marco teórico del proyecto.

\begin{description}[leftmargin=0pt, labelwidth=\textwidth, 
                    labelsep=0pt, style=nextline]

% -------------------------------------------------------
\item[\textbf{Adherencia al tratamiento}]
Grado de continuidad con que un paciente realiza las actividades de una 
intervención a lo largo del tiempo. En el contexto de la estimulación 
cognitiva domiciliaria, la adherencia es un factor crítico para la 
efectividad del programa, ya que los beneficios dependen de la práctica 
sostenida \parencite{moret, calatayud2021}.

% -------------------------------------------------------
\item[\textbf{Alucinación (LLMs)}]
Fenómeno propio de los Grandes Modelos de Lenguaje por el cual el modelo 
genera respuestas lingüísticamente plausibles pero factualmente incorrectas 
o inventadas. Representa una limitación relevante en sistemas de salud 
digital donde la precisión de la información es crítica \parencite{moret}.

% -------------------------------------------------------
\item[\textbf{API (\textit{Application Programming Interface})}]
Interfaz de programación de aplicaciones: conjunto de protocolos y 
definiciones que permite a un programa de software comunicarse con otro. 
En este proyecto se utiliza la Gemini API de Google para integrar el 
módulo de PLN al motor de juego Unity \parencite{Ortega2024}.

% -------------------------------------------------------
\item[\textbf{Arquitectura por capas}]
Patrón de diseño de software que organiza un sistema en módulos 
independientes con responsabilidades bien delimitadas que se comunican 
mediante interfaces estandarizadas. Se identifican tres capas recurrentes 
en sistemas de salud digital: presentación e interacción, inteligencia 
artificial y datos \parencite{Goumopoulos2021, Karapapas2021}.

% -------------------------------------------------------
\item[\textbf{Asistente virtual (VA)}]
Sistema de software basado en procesamiento de lenguaje natural capaz de 
interactuar con el usuario mediante lenguaje natural hablado o escrito. 
En herramientas de estimulación cognitiva, el asistente virtual cumple 
funciones de instrucción, motivación y retroalimentación adaptativa 
\parencite{moret}.

% -------------------------------------------------------
\item[\textbf{Atención sostenida}]
Capacidad cognitiva para mantener el foco en una tarea o estímulo durante 
un período prolongado sin pérdida de rendimiento. Es una de las funciones 
cognitivas más afectadas en el Deterioro Cognitivo Leve y constituye una 
esfera central de estimulación en este proyecto \parencite{bermejo2021, 
AbdAlrazaq2022}.

% -------------------------------------------------------
\item[\textbf{Design Science Research (DSR)}]
Paradigma metodológico de investigación aplicada cuyo objetivo es diseñar, 
construir y evaluar artefactos tecnológicos que resuelvan problemas del 
mundo real. A diferencia de la investigación descriptiva, el DSR produce 
como resultado principal el artefacto mismo y la evidencia de su 
funcionamiento \parencite{Munoz2025}.

% -------------------------------------------------------
\item[\textbf{Deterioro Cognitivo Leve (DCL)}]
Etapa clínica intermedia entre el envejecimiento cognitivo normal y la 
demencia, caracterizada por un declive perceptible en una o más funciones 
cognitivas —memoria, atención o funciones ejecutivas— que no compromete 
significativamente la autonomía funcional del individuo. Representa una 
ventana de oportunidad para la intervención preventiva 
\parencite{Giovagnoli2024}.

% -------------------------------------------------------
\item[\textbf{Dispositivo háptico}]
Hardware que transmite sensaciones físicas al usuario —resistencia, 
vibración o guía de movimiento— como canal de retroalimentación en 
sistemas interactivos. Se clasifican en dispositivos anclados al suelo 
(\textit{grounded}), de mano (\textit{hand-held}) y vestibles 
(\textit{wearable}) \parencite{Adilkhanov2022}.

% -------------------------------------------------------
\item[\textbf{DSR}]
Véase \textit{Design Science Research}.

% -------------------------------------------------------
\item[\textbf{Estimulación cognitiva (EC)}]
Conjunto de actividades y estrategias diseñadas para activar y ejercitar 
sistemáticamente las funciones cognitivas del individuo, con propósito 
preventivo y de mantenimiento —a diferencia de la rehabilitación clínica, 
que busca recuperar funciones perdidas—. Su eficacia ha sido documentada 
en memoria de trabajo, funciones ejecutivas y reducción de indicadores 
de depresión \parencite{bermejo2021, calatayud2021, arteaga2024}.

% -------------------------------------------------------
\item[\textbf{Exergame}]
Videojuego que requiere actividad física del usuario como mecanismo 
principal de interacción, combinando ejercicio motor con tareas cognitivas. 
En el contexto del envejecimiento activo, los exergames son una modalidad 
de entrenamiento dual con evidencia empírica de efectividad 
\parencite{exergames2026}.

% -------------------------------------------------------
\item[\textbf{Firebase}]
Plataforma de desarrollo de aplicaciones en la nube de Google, utilizada 
en este proyecto como infraestructura para integrar la Gemini API al motor 
de juego Unity mediante comunicación segura con la capa de inteligencia 
artificial del sistema \parencite{Ortega2024}.

% -------------------------------------------------------
\item[\textbf{Flexibilidad cognitiva}]
Capacidad para adaptar el pensamiento y la conducta ante cambios en las 
reglas o condiciones de una tarea. Es un componente central de las 
funciones ejecutivas y una de las habilidades más afectadas en el DCL 
\parencite{AbdAlrazaq2022, Giovagnoli2024}.

% -------------------------------------------------------
\item[\textbf{Force Feedback (FFB)}]
Mecanismo de retroalimentación háptica que utiliza motores de torque para 
ejercer resistencia activa sobre el eje de un dispositivo, simulando 
fuerzas físicas del entorno virtual. Se distingue de la vibración ERM en 
que genera resistencia real al movimiento del usuario 
\parencite{Marchal2008}.

% -------------------------------------------------------
\item[\textbf{Funciones ejecutivas (FE)}]
Conjunto de procesos cognitivos de alto nivel que permiten planificar, 
organizar, iniciar y supervisar conductas orientadas a un objetivo; 
comprenden memoria de trabajo, flexibilidad cognitiva, inhibición de 
respuestas y toma de decisiones. Son particularmente vulnerables al 
deterioro en adultos mayores \parencite{AbdAlrazaq2022}.

% -------------------------------------------------------
\item[\textbf{Gamificación}]
Aplicación de mecánicas y elementos propios del diseño de videojuegos 
—puntuación, progresión de niveles, retroalimentación inmediata, 
recompensas— en contextos no lúdicos, con el objetivo de incrementar la 
motivación y el compromiso del usuario \parencite{abd, Ortega2024}.

% -------------------------------------------------------
\item[\textbf{HID (\textit{Human Interface Device})}]
Clase de dispositivos de entrada y salida de datos estándar que se 
comunican con un ordenador sin requerir controladores adicionales, 
utilizando el protocolo USB HID. Los volantes comerciales como el 
Kontorland AP5 operan bajo este protocolo \parencite{Adilkhanov2022}.

% -------------------------------------------------------
\item[\textbf{Háptica}]
Rama de la interacción humano-computadora que estudia el sentido del tacto 
y la percepción cinestésica como canal de comunicación entre el usuario y 
el sistema. Permite transmitir información mediante sensaciones físicas 
como resistencia, vibración o presión \parencite{Marchal2008, 
Adilkhanov2022}.

% -------------------------------------------------------
\item[\textbf{IHC (Interacción Humano-Computadora)}]
Disciplina que estudia el diseño, evaluación e implementación de sistemas 
de cómputo para uso humano, con énfasis en la experiencia del usuario. 
Sus principios guían el diseño de interfaces accesibles para adultos 
mayores \parencite{Peres2022SUS}.

% -------------------------------------------------------
\item[\textbf{KPI (\textit{Key Performance Indicator})}]
Indicador clave de rendimiento: métrica cuantificable definida para 
evaluar el grado de cumplimiento de un objetivo. En el presente prototipo, 
los KPIs miden eficacia (tareas completadas correctamente), eficiencia 
(tiempo por tarea) y satisfacción (puntuación SUS) \parencite{Peres2022SUS, 
Ortega2024}.

% -------------------------------------------------------
\item[\textbf{LLM (\textit{Large Language Model})}]
Gran Modelo de Lenguaje: modelo de inteligencia artificial entrenado sobre 
enormes corpus de texto mediante arquitecturas de transformadores con 
mecanismos de auto-atención, capaz de generar texto contextualmente 
coherente y mantener conversaciones fluidas en lenguaje natural. Entre los 
más utilizados en salud digital se encuentran GPT (OpenAI), Gemini 
(Google) y LLaMA (Meta) \parencite{moret, Hamdi2025}.

% -------------------------------------------------------
\item[\textbf{MCI (\textit{Mild Cognitive Impairment})}]
Término en inglés equivalente al Deterioro Cognitivo Leve (DCL). Véase 
\textit{Deterioro Cognitivo Leve}.

% -------------------------------------------------------
\item[\textbf{Memoria de trabajo}]
Sistema cognitivo de capacidad limitada que permite mantener y manipular 
información de forma temporal durante la realización de tareas complejas. 
Es uno de los dominios cognitivos más frecuentemente afectados en el DCL 
y más entrenados en programas de estimulación cognitiva 
\parencite{AbdAlrazaq2022, bermejo2021}.

% -------------------------------------------------------
\item[\textbf{MoCA (\textit{Montreal Cognitive Assessment})}]
Instrumento de tamizaje cognitivo de 30 puntos, diseñado específicamente 
para detectar el Deterioro Cognitivo Leve con mayor sensibilidad (87\%) 
que el MMSE. Evalúa memoria, atención, funciones ejecutivas, lenguaje, 
habilidades visuoespaciales y orientación. Existe versión validada en 
español \parencite{Giovagnoli2024}.

% -------------------------------------------------------
\item[\textbf{MMSE (\textit{Mini-Mental State Examination})}]
Instrumento de tamizaje cognitivo de 30 puntos ampliamente utilizado en 
la clínica. Presenta menor sensibilidad para DCL (72--94\%) que el MoCA 
y efecto techo en personas con alta escolaridad. Su acceso requiere 
licencia \parencite{Giovagnoli2024}.

% -------------------------------------------------------
\item[\textbf{Motor ERM (\textit{Eccentric Rotating Mass})}]
Motor de vibración de masa excéntrica rotatoria: actuador que genera 
vibración mediante el giro de una masa desequilibrada. Es el mecanismo 
háptico predominante en controladores comerciales de bajo costo como 
mandos de consola y volantes de entrada. Genera patrones de vibración 
programables sin ejercer resistencia activa sobre el eje del dispositivo 
\parencite{Adilkhanov2022, Gardinger2024}.

% -------------------------------------------------------
\item[\textbf{Motor de juego (\textit{game engine})}]
Plataforma de software que proporciona los componentes fundamentales para 
el desarrollo de aplicaciones interactivas: renderizado gráfico, gestión 
de físicas, manejo de entradas, audio y herramientas de scripting. En este 
proyecto se utiliza Unity como motor principal \parencite{Ortega2024}.

% -------------------------------------------------------
\item[\textbf{Orientación espacial}]
Capacidad cognitiva para ubicarse y desplazarse en el entorno, 
reconociendo la posición propia en relación con el espacio circundante. 
Es uno de los dominios cognitivos objeto de estimulación en este prototipo, 
trabajado mediante la navegación en laberintos \parencite{bermejo2021}.

% -------------------------------------------------------
\item[\textbf{PLN (Procesamiento de Lenguaje Natural)}]
Rama de la inteligencia artificial que estudia la interacción entre 
computadoras y el lenguaje humano, habilitando capacidades como 
comprensión, generación y traducción de texto o voz de forma automatizada. 
Es la base tecnológica del asistente conversacional de este prototipo 
\parencite{moret}.

% -------------------------------------------------------
\item[\textbf{Praxias visomotoras}]
Capacidad para ejecutar movimientos voluntarios aprendidos en respuesta 
a estímulos visuales, requiriendo coordinación entre la percepción visual 
y la respuesta motora. En este proyecto se activan mediante la coordinación 
simultánea de manos y pies durante la interacción con el volante y el pedal 
\parencite{bermejo2021, AbdAlrazaq2022}.

% -------------------------------------------------------
\item[\textbf{\textit{Prompt}}]
Instrucción o texto de entrada que se proporciona a un LLM para guiar su 
respuesta. En sistemas de salud digital, el uso de prompts estructurados 
y acotados es una estrategia para mitigar el fenómeno de alucinación y 
mantener la interacción dentro de contextos controlados y seguros 
\parencite{Ortega2024, moret}.

% -------------------------------------------------------
\item[\textbf{Privacidad por diseño (\textit{Privacy by Design})}]
Principio de arquitectura de software que integra la protección de datos 
personales desde las etapas iniciales del diseño del sistema, en lugar 
de añadirla como medida posterior. En sistemas de salud domiciliaria 
implica almacenar datos localmente y minimizar la transmisión de 
información identificable a servicios en la nube \parencite{moret, 
Munoz2025}.

% -------------------------------------------------------
\item[\textbf{Reserva cognitiva (RC)}]
Capacidad del cerebro para tolerar los efectos de la patología neurológica 
sin manifestar síntomas clínicos, actuando como mecanismo de compensación 
activa. Se forja a través de la educación, la actividad intelectual y el 
compromiso social a lo largo de la vida \parencite{Stern2022, 
Giovagnoli2024}.

% -------------------------------------------------------
\item[\textbf{Serious Game}]
Aplicación de software que utiliza mecánicas y elementos propios del 
diseño de videojuegos con un propósito primario distinto al entretenimiento: 
en el contexto de la salud, su objetivo es promover comportamientos 
terapéuticos, educativos o de rehabilitación, incorporando retroalimentación 
inmediata, progresión de dificultad y motivación intrínseca 
\parencite{abd, Ortega2024}.

% -------------------------------------------------------
\item[\textbf{SUS (\textit{System Usability Scale})}]
Escala de usabilidad percibida de diez ítems que genera una puntuación 
de 0 a 100. Una puntuación $\geq$ 68 indica usabilidad aceptable; 
puntuaciones $>$ 80.3 se clasifican como excelentes para aplicaciones de 
salud digital. Es el instrumento de evaluación de usabilidad seleccionado 
para la prueba piloto de este proyecto \parencite{Peres2022SUS}.

% -------------------------------------------------------
\item[\textbf{Tarea dual (\textit{dual-task})}]
Paradigma experimental y de entrenamiento cognitivo que exige la 
ejecución simultánea de dos tareas —una motriz y una cognitiva—, 
revelando déficits atencionales que permanecen ocultos bajo condiciones 
de tarea simple. El entrenamiento dual ha demostrado mejoras 
significativamente mayores en funciones ejecutivas que el entrenamiento 
cognitivo aislado \parencite{Voelcker2025DualTask, Li2023DualTask}.

% -------------------------------------------------------
\item[\textbf{Tareas ecológicamente válidas}]
Actividades de entrenamiento o evaluación cognitiva que simulan 
situaciones de la vida cotidiana reconocibles por el usuario. En el diseño 
de Serious Games para adultos mayores, la validez ecológica favorece la 
adherencia y la transferencia de las habilidades entrenadas al contexto 
real \parencite{Ortega2024, Gardinger2024}.

% -------------------------------------------------------
\item[\textbf{Unity}]
Motor de juego multiplataforma desarrollado por Unity Technologies, con 
soporte nativo para dispositivos HID, scripting en C\# y conexión con 
APIs externas. Es el motor de mayor adopción en proyectos de Serious Games 
de salud y rehabilitación cognitiva \parencite{Ortega2024, Ribeiro2024}.

% -------------------------------------------------------
\item[\textbf{Usabilidad}]
Grado en que un sistema permite a usuarios específicos alcanzar objetivos 
concretos con efectividad, eficiencia y satisfacción en un contexto de uso 
determinado. En herramientas para adultos mayores, la usabilidad es un 
criterio de diseño prioritario dada la menor experiencia de este grupo con 
interfaces digitales \parencite{Peres2022SUS, Ortega2024}.

% -------------------------------------------------------
\item[\textbf{Volante háptico de bajo costo}]
Controlador de simulación de conducción de precio accesible (aproximadamente 
\$20--\$40 USD) que integra motores ERM para retroalimentación vibrotáctil 
y se comunica mediante protocolos HID estándar (XInput/DInput). El modelo 
Kontorland AP5, utilizado en este prototipo, pertenece a esta categoría 
\parencite{Adilkhanov2022, Gardinger2024}.

% -------------------------------------------------------
\item[\textbf{XInput / DInput}]
Protocolos estándar de Windows para la comunicación entre controladores 
de juego y aplicaciones de software. XInput es el estándar moderno para 
controles Xbox-compatibles; DInput (DirectInput) es el protocolo heredado 
que soporta una gama más amplia de dispositivos HID, incluyendo volantes 
comerciales de bajo costo \parencite{Adilkhanov2022}.

\end{description}
 %   \cleardoublepage

    % --- INICIO DEL CONTENIDO PRINCIPAL ---
    \pagenumbering{arabic} 
    \setcounter{page}{1}   

    % =====================================================================
%  \nuevo{...} resalta en amarillo el texto de la versión 2.
%  Para la versión sin resaltado: \resaltarfalse en main.tex.
% =====================================================================
\chapter[Protocolo de investigación]{\nuevo{Protocolo de investigación}}

\section{Antecedentes}

La búsqueda de vivienda en renta por parte de estudiantes universitarios que provienen de otras localidades distintas a Miahuatlán de Porfirio Díaz es un problema que ha sido abordado desde distintas perspectivas: el estudio del impacto de las universidades en el desarrollo local, el análisis del crecimiento urbano, el efecto de las plataformas digitales en el mercado de alquiler\nuevo{, la confianza en las transacciones en línea entre particulares, los mecanismos de verificación, el fraude en la renta en línea} y el desarrollo de sistemas informáticos para la gestión de arrendamientos. En esta sección se presenta, en primer lugar, el panorama estadístico y territorial que enmarca el problema y, posteriormente, una síntesis crítica de las investigaciones previas más relevantes, con el propósito de identificar lo que ya se conoce, las limitaciones de las propuestas existentes y la brecha que el presente proyecto busca atender.

\subsection{Panorama estadístico y territorial}

De acuerdo con la Encuesta Nacional de Vivienda (ENVI) 2020, levantada mediante un muestreo probabilístico, bietápico, estratificado y por conglomerados de 55\,147 viviendas, de los 35.3 millones de viviendas particulares habitadas en México, el 16.4\,\% son rentadas; de estas 5.8 millones de viviendas en renta, únicamente el 54.0\,\% cuenta con un contrato de renta vigente \parencite{INEGI2021ENVI}. Es decir, casi la mitad de los arrendamientos del país se realiza sin un instrumento formal que dé certeza a arrendadores y arrendatarios. Asimismo, de acuerdo con el \textcite{INEGI2021ENVI}, el 21.1\,\% de los hogares (7.6 millones) reportó que alguno de sus integrantes necesita o planea rentar, comprar o construir una vivienda. Aunque la ENVI tiene representatividad nacional y estatal, no ofrece información a nivel municipal ni distingue la demanda específica de la población estudiantil.

En el ámbito local, \nuevo{los datos del Censo de Población y Vivienda 2020 que difunde la \textcite{SE2020DataMexico}} a través de la plataforma Data México \nuevo{describen} el perfil habitacional de Miahuatlán de Porfirio Díaz.
\nuevo{El municipio contaba en 2020 con 50\,375 habitantes (51.2\,\% mujeres y 48.8\,\% hombres) y con alrededor de 13\,100 viviendas, lo que equivale a cerca de 3.8 habitantes por vivienda; la mayoría de las viviendas particulares habitadas tenía tres cuartos (28.2\,\%) o dos cuartos (23.4\,\%). La plataforma no reporta, a nivel municipal, la proporción de viviendas rentadas.}
Estos datos corresponden al corte censal de 2020 y no reflejan la oferta privada de renta ni la ubicación de los inmuebles disponibles. De forma complementaria, el Plan de Ordenamiento Territorial elaborado por el \textcite{SEINFRA2024} delimita las colonias del municipio, la cobertura de las redes de servicios básicos y la proximidad a nodos educativos como la Universidad de la Sierra Sur (UNSIS), con base en un diagnóstico territorial participativo y análisis cartográfico. Al ser un instrumento de planeación estático, no permite consultar la oferta de vivienda en renta, pero proporciona la base geográfica oficial para categorizar las zonas en las que se ubican los inmuebles.

\subsection{La universidad como generadora de demanda habitacional en Miahuatlán de Porfirio Díaz}

\textcite{ZepedaGarcia2019} examinaron cómo la llegada de la UNSIS ha transformado a Miahuatlán de Porfirio Díaz en un periodo aproximado de quince años, con el objetivo de valorar a la universidad como productora de capital humano y como impulsora de la economía local. Mediante una investigación de campo que recuperó testimonios de habitantes que presenciaron dichos cambios, complementada con análisis documental, encontraron que la llegada de estudiantes, profesores y personal administrativo modificó las condiciones de vida de la ciudad y generó una demanda creciente de servicios, entre ellos el alojamiento. Sus resultados, sin embargo, son de carácter cualitativo: no cuantifican la demanda de vivienda ni proponen mecanismos para organizarla. Este trabajo es relevante para el presente proyecto porque documenta, en el mismo contexto geográfico, que la comunidad universitaria constituye \nuevo{una fuente} de demanda de vivienda en renta.

En una línea complementaria, \textcite{AvilaFlores2020} analizaron los factores del crecimiento urbano y del desarrollo sustentable en Miahuatlán de Porfirio Díaz durante el siglo \textsc{xxi}. A través de un análisis territorial, espacial y demográfico apoyado en Sistemas de Información Geográfica (SIG) y cartografía, estimaron la tasa de expansión urbana y la densidad de población, lo que les permitió identificar las zonas de mayor presión y expansión dentro del municipio. Su enfoque es macroscópico, por lo que no llega al nivel de la oferta inmobiliaria individual. Para esta investigación, dicho estudio aporta elementos para fundamentar el componente de geolocalización de la plataforma, al permitir identificar las zonas de expansión en las que resulta más necesario vincular la oferta de cuartos con los estudiantes.
% TODO: incorporar una cifra concreta de este estudio (tasa de expansión urbana o densidad) tras revisar el texto completo.

\subsection{Las plataformas digitales en el mercado de vivienda en alquiler}

\nuevo{\textcite{GarridoRodriguez2025} analizó la distribución geográfica y los precios de los alojamientos ofertados en Airbnb en la Ciudad de México entre 2019 y 2023, con datos de InsideAirbnb y del Censo de Población y Vivienda. Encontró que la oferta se recuperó tras la pandemia hasta superar la de 2019 y que en algunos territorios más del 60\,\% del parque habitacional se encuentra inserto en la plataforma. Su análisis se limita al alquiler turístico de corta estancia en una metrópoli, por lo que su pertinencia para este proyecto es indirecta: sugiere que las plataformas digitales reorganizan el mercado de alquiler y que el análisis geográfico de la oferta es útil para comprenderlo, aunque no se identificaron estudios equivalentes para el alquiler de mediana estancia dirigido a estudiantes en municipios pequeños.}

\subsection[Confianza y verificación en plataformas entre particulares]{\nuevo{Confianza y verificación en plataformas entre particulares}}

\nuevo{La confianza percibida es una de las dos variables que contrasta la hipótesis de este proyecto, por lo que se revisaron los trabajos que la estudian en el comercio en línea y en las plataformas entre particulares. \textcite{McKnight2002} propusieron y validaron empíricamente un conjunto de medidas de confianza para el comercio electrónico que distingue la disposición a confiar, la confianza basada en la institución (la percepción de que existen estructuras y garantías que protegen al usuario) y las creencias de confianza sobre el vendedor (competencia, benevolencia e integridad). Esta tipología es relevante porque separa la confianza en el medio de la confianza en la contraparte, que en el caso de la renta de vivienda es el arrendador. \textcite{Gefen2003} integraron la confianza en el modelo de aceptación de tecnología y mostraron, con compradores en línea, que la confianza influye en la intención de uso junto con la utilidad y la facilidad de uso percibidas, y que se apoya en la percepción de garantías estructurales del sitio. En el ámbito de la economía colaborativa, \textcite{Ert2016} mostraron que en Airbnb las señales visuales sobre el anfitrión, como su fotografía, se asocian con el precio y la probabilidad de elegir un alojamiento, mientras que las calificaciones de reputación tuvieron un efecto menor, lo que sugiere que los elementos que el usuario interpreta como indicios de la identidad y honestidad del oferente pesan en su decisión.}

\nuevo{La confianza en el medio, que es la que este proyecto contrasta, corresponde a lo que \textcite{Pavlou2004} denominan confianza basada en la institución: la percepción del comprador de que existen mecanismos institucionales de terceros, como la retroalimentación, las garantías o la intermediación, que favorecen el éxito de la transacción. En un estudio con compradores de un mercado en línea, estos autores mostraron que la efectividad percibida de dichos mecanismos genera confianza no solo en unos pocos vendedores reconocidos, sino en la comunidad de vendedores en su conjunto, y que esa confianza reduce el riesgo percibido. Este resultado sostiene la idea de que un medio que incorpora mecanismos de verificación puede elevar la confianza del estudiante en las ofertas que publica, aun cuando no conozca a cada arrendador.}

\nuevo{Los mecanismos de verificación y reputación cuentan con literatura propia. \textcite{Resnick2000} describieron los sistemas de reputación como mecanismos que recopilan y difunden la retroalimentación sobre el comportamiento pasado de los participantes, para fomentar la confianza entre desconocidos en internet, y señalaron sus límites, entre ellos la dificultad de obtener retroalimentación, en especial la negativa. \textcite{Ba2002} mostraron, con compradores de subastas en línea, que
 las calificaciones de los vendedores generan confianza en su credibilidad y benevolencia, y que esta confianza se traduce en primas de precio. En el ámbito del alojamiento entre particulares, \textcite{Ma2017} analizaron cerca de 1200 perfiles de anfitriones de Airbnb y encontraron, mediante anotación de contenidos y un experimento, que la información que el anfitrión revela sobre sí mismo se relaciona con la confiabilidad que perciben los huéspedes. \textcite{Zervas2021} compararon las calificaciones de los mismos alojamientos en Airbnb y en TripAdvisor y hallaron que en Airbnb casi todas son altas, lo que reduce su capacidad para distinguir entre ofertas. Para el presente proyecto, estos estudios sugieren que las reseñas por sí solas son un mecanismo insuficiente en un mercado pequeño y con pocas transacciones registradas, y que conviene complementarlas con una verificación directa de la identidad del arrendador y de la existencia del inmueble.}

\subsection[Fraude en la renta de vivienda en línea]{\nuevo{Fraude en la renta de vivienda en línea}}

\nuevo{\textcite{Park2017} realizaron el primer estudio empírico sistemático sobre fraudes en anuncios de renta en línea. Analizaron más de dos millones de anuncios publicados en Craigslist durante cinco meses en 20 ciudades de Estados Unidos, detectaron alrededor de 29\,000 anuncios fraudulentos y, mediante un sistema automatizado que conversaba con los estafadores, identificaron siete tipos de fraude, entre ellos la clonación de anuncios reales para cobrar un depósito por adelantado. El sitio eliminó solo cerca de la mitad de esos anuncios, en promedio unas diez horas después de su publicación. Aunque el estudio se ubica en un mercado distinto del mexicano, documenta que los anuncios no verificados facilitan el fraude y que la moderación posterior a la publicación es insuficiente. En México, la Secretaría de Seguridad y Protección Ciudadana alertó en 2025 sobre anuncios de renta de vivienda publicados en redes sociales con precios inusualmente bajos, en los que se usan identificaciones falsas y contratos apócrifos para solicitar depósitos por adelantado, tras lo cual los anuncios desaparecen \parencite{AnimalPolitico2025}. No se encontraron registros de este tipo de fraude para Miahuatlán de Porfirio Díaz, por lo que su frecuencia local se documentará como parte del objetivo específico 1.}

\subsection{Sistemas web para la gestión de arrendamientos}

El antecedente tecnológico más cercano a esta investigación es el de \textcite{MejiaAvecias2023}, quienes desarrollaron un sistema web para facilitar la obtención y distribución de información sobre departamentos en renta para estudiantes foráneos de la Escuela Superior de Tlahuelilpan, Hidalgo. A partir de información recabada con estudiantes de la zona, identificaron que la oferta se difundía principalmente por redes sociales, con información escasa, ambigua y dispersa, e incluso con publicaciones de más de un año de antigüedad sin actualizar. El desarrollo siguió la metodología ágil Scrum; los requerimientos se obtuvieron mediante entrevistas a estudiantes foráneos y el sistema se construyó con una arquitectura cliente-servidor empleando HTML, CSS, JavaScript, PHP y MySQL. Como resultado, presentaron el modelo de datos, el diagrama de módulos y doce interfaces que incluyen el registro de inmuebles, la búsqueda con filtros por ubicación, tamaño, servicios y costo, la reservación en línea y la comunicación entre arrendadores y estudiantes. Asimismo, al revisar plataformas existentes como Stukey, Roomgo o Spotahome, señalan que estas se concentran en la Ciudad de México y en otras ciudades grandes del país. No obstante, el trabajo no reporta una evaluación empírica del sistema, como pruebas con usuarios, métricas de usabilidad o tiempos de búsqueda, y fue diseñado para el contexto de Tlahuelilpan. Este antecedente sirve como guía para definir el modelo de datos, los módulos y los filtros de búsqueda de la plataforma propuesta, la cual, a diferencia de este, contemplará la evaluación con usuarios y las condiciones propias de Miahuatlán de Porfirio Díaz.

\nuevo{En un contexto similar al de Miahuatlán, \textcite{Bisquera2026} desarrollaron y evaluaron un localizador de casas de huéspedes basado en sistemas de información geográfica para los estudiantes de la Nueva Vizcaya State University, en un municipio de Filipinas. La plataforma integra un mapa interactivo con rutas, filtros por precio y ubicación, distancia al campus y actualización de la disponibilidad por parte de los propietarios, además del registro de los permisos de operación de cada inmueble. Al compararla con los métodos tradicionales (recorridos a pie y recomendaciones de conocidos), reportaron un ahorro de 181.25 minutos por búsqueda y de 35 pesos filipinos en traslados, y en una prueba de usabilidad con 175 participantes obtuvieron una calificación media de 4.83 sobre 5. La comparación, sin embargo, se basa en estimaciones de los propios usuarios y no en un diseño experimental, y la plataforma no mide la confianza. Por su parte, \textcite{Oluwadare2025} combinaron entrevistas con residentes y administradores de albergues estudiantiles con herramientas de SIG para construir una base de datos y mapas de la oferta de alojamiento alrededor de la Obafemi Awolowo University, en Nigeria; encontraron que los albergues se concentran en los límites del campus y que la elección depende del costo, la disponibilidad de electricidad, la seguridad, los servicios y la cercanía, y recomendaron mantener actualizada la base de datos mediante una plataforma web.}

\nuevo{Entre las plataformas comerciales, \textcite{MejiaAvecias2023} señalan que Stukey, Roomgo y Spotahome se concentran en la Ciudad de México y en otras ciudades grandes. De ellas, Spotahome es la que describe un mecanismo de verificación explícito: verifica a cada arrendador antes de publicar su inmueble y, cuando es posible, envía personal propio a visitar, filmar y medir la vivienda, de modo que cada anuncio incluye un recorrido en video, fotografías y un plano \parencite{Spotahome2026}. Opera en ciudades europeas, por lo que no atiende mercados como el de Miahuatlán. Este modelo muestra que la verificación presencial es factible, pero también que depende de personal dedicado, lo que en un municipio pequeño obliga a diseñar un mecanismo más ligero.}

\nuevo{De forma más lejana,} \textcite{Valenzuela2024} desarrolló un sistema web para el control de pagos por arrendamiento de locales \nuevo{comerciales} de la Municipalidad Provincial de Pasco, Perú, \nuevo{y reportó mejoras en la rapidez y confiabilidad del registro de pagos con una muestra de 20 arrendatarios. Aunque su objeto (pagos de locales municipales) es distinto del alojamiento estudiantil, se retoma únicamente como referencia para el modelado de datos de un sistema de rentas.}

\subsection{Síntesis y brecha identificada}
Los estudios sobre Miahuatlán de Porfirio Díaz \parencite{AvilaFlores2020, ZepedaGarcia2019} \nuevo{documentan} que la presencia de la \nuevo{UNSIS} genera demanda \nuevo{de vivienda} y que existen herramientas de análisis territorial para ubicarla; sin embargo, ninguno propone una solución que organice la oferta. La literatura sobre plataformas digitales \parencite{GarridoRodriguez2025} \nuevo{sugiere} que estas reorganizan el mercado de alquiler, pero se ha concentrado en el alquiler turístico de las grandes ciudades. \nuevo{Los estudios sobre confianza en línea \parencite{McKnight2002, Gefen2003, Ert2016} muestran que la confianza en el medio y en la contraparte influye en la decisión del usuario, y que los mecanismos institucionales del medio pueden generar confianza en el conjunto de oferentes \parencite{Pavlou2004}; a su vez, la literatura sobre reputación \parencite{Resnick2000, Ba2002, Ma2017, Zervas2021} sugiere que las reseñas por sí solas son insuficientes, y los estudios sobre fraude \parencite{Park2017} documentan el riesgo de los anuncios no verificados. No se identificaron trabajos que midan la confianza en la búsqueda de vivienda estudiantil.} Finalmente, los sistemas web existentes \parencite{MejiaAvecias2023, Valenzuela2024, Bisquera2026, Oluwadare2025} \nuevo{muestran que es técnicamente posible} centralizar la información de arrendamiento, pero fueron diseñados para otros contextos o tipos de inmueble y no reportan una evaluación \nuevo{experimental} con usuarios finales\nuevo{; el que más se acerca \parencite{Bisquera2026} compara tiempos estimados por los usuarios y no incluye la confianza ni un mecanismo de verificación de arrendadores}.

En consecuencia, no se identificó una solución que integre la oferta de cuartos y departamentos en renta de Miahuatlán de Porfirio Díaz con información geográfica sobre colonias, servicios y cercanía a la UNSIS, \nuevo{que incorpore mecanismos de verificación de arrendadores e inmuebles,} dirigida a \nuevo{estudiantes} no residentes y evaluada con sus usuarios. Esta es la brecha que el presente proyecto busca atender.

\section{Planteamiento del problema}

Miahuatlán de Porfirio Díaz, Oaxaca, se ha consolidado progresivamente como un importante centro comercial y académico de la región, lo que ha propiciado la llegada de personas provenientes de distintas localidades. Entre este sector de la población destacan los estudiantes de la \nuevo{UNSIS}, cuya matrícula es de \nuevo{2294 estudiantes \parencite{UNSIS2026ServiciosEscolares},}
% PENDIENTE (en gestión con Servicios Escolares): confirmar la cifra y el ciclo escolar, y obtener el porcentaje de estudiantes foráneos.
de los cuales \nuevo{el [\textit{porcentaje}]\,\% proviene de otras localidades \parencite{UNSIS2026ServiciosEscolares} y}
requiere trasladarse temporalmente a Miahuatlán para continuar con su formación académica. Esta situación genera una demanda constante de viviendas en renta entre quienes no cuentan con un lugar de residencia permanente en la localidad.

Actualmente, la difusión de la oferta de viviendas en renta se realiza principalmente mediante grupos de redes sociales, como Facebook, así como a través de anuncios físicos colocados en postes y espacios cercanos a las instalaciones de la universidad.
\nuevo{Para caracterizar estos medios se realizará un registro exploratorio de los grupos de Facebook dedicados a la renta en el municipio (número de grupos activos, publicaciones por mes y antigüedad de las publicaciones) y de los anuncios físicos colocados cerca de la universidad.} Estos medios funcionan de manera independiente y no existe un canal centralizado que permita consultar, organizar y actualizar la información de los inmuebles disponibles. Esta dispersión dificulta el acceso a información confiable y oportuna, ya que las publicaciones pueden quedar desactualizadas, presentar datos incompletos o carecer de elementos que permitan verificar la identidad de los arrendadores y las condiciones reales de las viviendas ofertadas.

Como consecuencia, \nuevo{los estudiantes no residentes} enfrentan dificultades para localizar y comparar alternativas de vivienda de acuerdo con su presupuesto, ubicación, características y necesidades particulares. Esto incrementa el tiempo y esfuerzo dedicado al proceso de búsqueda, además de generar incertidumbre al momento de establecer contacto con los arrendadores. Asimismo, la falta de mecanismos de verificación \nuevo{expone a los estudiantes al riesgo de fraudes}, como ofertas de inmuebles inexistentes, publicaciones engañosas o viviendas cuyas condiciones no corresponden con lo anunciado.
\nuevo{No existe un registro local de estos casos; a nivel nacional, las autoridades han alertado sobre anuncios de renta falsos en redes sociales que se usan para cobrar depósitos por adelantado \parencite{AnimalPolitico2025}. Para dimensionar el problema en Miahuatlán se aplicará un sondeo exploratorio a estudiantes no residentes de la UNSIS sobre su experiencia con ofertas falsas o con condiciones distintas a las anunciadas.}

\nuevo{El problema tiene también un lado de oferta. Los arrendadores locales dependen de publicaciones que pierden visibilidad en pocos días dentro de los grupos de redes sociales y de anuncios físicos de alcance limitado, por lo que deben volver a publicar con frecuencia y reciben consultas de personas cuyas necesidades no corresponden con su inmueble. Una plataforma centralizada solo resulta útil si los arrendadores publican en ella; por ello, este proyecto considera también sus necesidades y las condiciones bajo las cuales estarían dispuestos a adoptar un canal nuevo y a someterse a un proceso de verificación.}
% TODO: confirmar esta descripción con las entrevistas a arrendadores (objetivo específico 1).

En este contexto, se identifica la necesidad de contar con una solución tecnológica que centralice la oferta de viviendas en renta en Miahuatlán de Porfirio Díaz, facilite la búsqueda y comparación de inmuebles mediante criterios específicos\nuevo{,} proporcione información organizada y actualizada \nuevo{e incorpore mecanismos de verificación que contribuyan a generar confianza entre arrendadores y estudiantes}.

\section[Preguntas de investigación]{\nuevo{Preguntas de investigación}}

\subsection[Pregunta general]{\nuevo{Pregunta general}}

\nuevo{¿En qué medida el uso de una plataforma web centralizada, con geolocalización y verificación de inmuebles, reduce el tiempo de localización de ofertas y aumenta la confianza percibida de los estudiantes no residentes de la Universidad de la Sierra Sur, en comparación con los medios de difusión actuales?}

\subsection[Preguntas específicas]{\nuevo{Preguntas específicas}}

\begin{enumerate}
    \item \nuevo{¿Qué necesidades, criterios de búsqueda y riesgos reportan los estudiantes no residentes y los arrendadores?}
    \item \nuevo{¿Qué mecanismos de verificación son viables y aceptables para arrendadores y estudiantes?}
    \item \nuevo{¿Difieren el tiempo de localización de ofertas, la exactitud de la elección y la confianza percibida entre la plataforma y los medios actuales?}
    \item \nuevo{¿Qué nivel de usabilidad percibida alcanza el prototipo?}
\end{enumerate}

\section{Justificación}

El presente proyecto se justifica en cuatro \nuevo{dimensiones (social, tecnológica, económica y metodológica) y se complementa con el análisis de su viabilidad y continuidad.}

\textbf{Social.} Los beneficiarios directos son los estudiantes no residentes de la UNSIS, \nuevo{que representan el [\textit{porcentaje}]\,\% de una matrícula de 2294 estudiantes \parencite{UNSIS2026ServiciosEscolares}}. Para esta población, la búsqueda de vivienda depende hoy de publicaciones dispersas en redes sociales y de anuncios físicos que no permiten verificar la existencia del inmueble ni la identidad del arrendador. Una plataforma que concentre la oferta y que incorpore mecanismos de verificación \nuevo{puede contribuir} a reducir la exposición a ofertas engañosas o inexistentes y a que los estudiantes elijan su alojamiento con información suficiente sobre su ubicación, servicios y condiciones, lo que favorece su adaptación al municipio durante su formación académica.

\textbf{Tecnológica.} Los sistemas web reportados en la literatura para la renta de vivienda estudiantil \parencite{MejiaAvecias2023} \nuevo{muestran que es técnicamente posible} centralizar esta información, pero no integran componentes geográficos ni reportan una evaluación con usuarios finales. El proyecto aporta una plataforma que vincula cada inmueble con la delimitación oficial de colonias y la cobertura de servicios del Plan de Ordenamiento Territorial \parencite{SEINFRA2024}, así como con su distancia a la UNSIS, lo que permite filtrar y comparar opciones con criterios que \nuevo{no se encontraron en los medios de difusión locales revisados}.

\textbf{Económica.} Para los arrendadores locales, la plataforma representa un canal de difusión de bajo costo que amplía la visibilidad de sus inmuebles hacia la comunidad universitaria, sin depender de la vigencia de una publicación en grupos de redes sociales. Para los \nuevo{estudiantes}, la posibilidad de comparar precios, servicios y ubicación en un solo lugar facilita la elección de la opción con mejor relación costo--beneficio para su presupuesto\nuevo{. Como beneficio esperado, que no será medido en este estudio, se prevé una reducción de} los gastos asociados a traslados y visitas a inmuebles que no cumplen lo anunciado. Asimismo, el registro estructurado de las condiciones de renta \nuevo{(precio, depósito, servicios incluidos y reglas) ofrece a ambas partes una referencia clara al momento de acordar la renta}, en un país donde casi la mitad de las viviendas rentadas carece de contrato vigente \parencite{INEGI2021ENVI}. \nuevo{La plataforma no genera contratos ni sustituye la asesoría legal.}

\textbf{Metodológica.} A diferencia de los antecedentes revisados, el proyecto contempla la evaluación de la plataforma con sus usuarios mediante \nuevo{un diseño experimental e} indicadores medibles\nuevo{: el tiempo de localización de ofertas, la exactitud de la elección, la confianza percibida y la usabilidad}. El protocolo de evaluación diseñado podrá reutilizarse en otros municipios con presencia universitaria y un mercado de renta informal similar al de Miahuatlán de Porfirio Díaz, para el cual no se identificaron estudios previos.

\nuevo{\textbf{Viabilidad.} El proyecto se desarrollará durante el periodo 2026--2027 con los recursos disponibles para las tesistas: equipo de cómputo propio, herramientas de desarrollo de código abierto y servicios de alojamiento web de bajo costo.}
\nuevo{La visualización geográfica se apoyará en los mapas de OpenStreetMap, cuyos datos son abiertos (licencia ODbL), mediante la biblioteca de código abierto Leaflet, por lo que no implica costo de licencia, a diferencia de la API de Google Maps, cuyo uso se cobra a partir de cierto volumen de consultas. La cartografía de colonias proviene del Plan de Ordenamiento Territorial \parencite{SEINFRA2024}; si no se obtiene en formato vectorial (shapefile o GeoJSON), los límites de las colonias se digitalizarán sobre OpenStreetMap a partir de los planos publicados en el Plan. El acceso a los participantes se facilita porque las tesistas forman parte de la comunidad de la UNSIS; el acceso a los arrendadores se gestionará a partir de las ofertas publicadas en los medios actuales, solicitando su autorización para incluir sus inmuebles en la plataforma.}

\nuevo{\textbf{Continuidad.} Al concluir la tesis, la plataforma requiere que alguien mantenga vigente la oferta y atienda las verificaciones.}
% PENDIENTE (en gestión): confirmar con el Instituto de Informática que aceptará la operación de la plataforma.
\nuevo{Se propone transferir su operación al Instituto de Informática de la UNSIS, para que la administren estudiantes de servicio social o de proyectos de la licenciatura bajo la supervisión de un profesor, junto con la documentación técnica y un manual de administración. Además, la plataforma se diseñará para reducir la carga de mantenimiento: los arrendadores actualizarán la disponibilidad de sus inmuebles y las publicaciones que no se confirmen en 30 días dejarán de mostrarse automáticamente.}

\section{Objetivos}

\subsection{Objetivo general}

\nuevo{Desarrollar} una plataforma web centralizada, con geolocalización y verificación de inmuebles, \nuevo{y evaluar su efecto} en el tiempo de \nuevo{localización de ofertas} y en la confianza percibida \nuevo{de} estudiantes no residentes de la Universidad de la Sierra Sur en Miahuatlán de Porfirio Díaz, Oaxaca, durante el periodo 2026--2027.

\subsection{Objetivos específicos}

\begin{enumerate}
    \item Identificar las necesidades, criterios de búsqueda y riesgos \nuevo{de} estudiantes no residentes y arrendadores en el proceso de renta.
    \item Diseñar la arquitectura, el modelo de datos y los mecanismos de verificación de arrendadores e inmuebles de la plataforma, con base en los requerimientos identificados y en la delimitación territorial de colonias y servicios del municipio.
    \item Implementar un prototipo funcional que integre registro de inmuebles, \nuevo{verificación de arrendadores e inmuebles,} búsqueda por criterios y visualización geográfica.
    \item \nuevo{Construir el conjunto de ofertas reales que alimentará ambas condiciones de la evaluación: recopilarlas, obtener la autorización de los arrendadores, geolocalizarlas y verificarlas.}
    \item Comparar el tiempo de \nuevo{localización de ofertas, la exactitud de la elección y la} confianza percibida de un grupo de estudiantes al \nuevo{buscar} vivienda mediante la plataforma frente a los medios de difusión actuales.
    \item Evaluar la usabilidad \nuevo{percibida} del prototipo con estudiantes y arrendadores mediante la \nuevo{escala SUS}.
\end{enumerate}

\section{Hipótesis}

La presente investigación tiene un enfoque cuantitativo y \nuevo{emplea un diseño experimental intrasujetos con contrabalanceo: cada participante realiza la búsqueda con ambos medios, y el orden de uso se asigna al azar} \parencite{HernandezSampieri2018}.

\subsection{Hipótesis de investigación}

El uso de una plataforma web centralizada, con geolocalización y mecanismos de verificación de inmuebles, reduce el tiempo de \nuevo{localización de ofertas de vivienda} en renta y aumenta \nuevo{la} confianza percibida \nuevo{en el medio} de los estudiantes no residentes de la UNSIS en comparación con los medios de difusión actuales (grupos de Facebook y anuncios físicos). Elementos \nuevo{de} la hipótesis:

\begin{itemize}
    \item \textbf{Unidad de análisis:} estudiantes no residentes de la Universidad de la Sierra Sur que buscan o han buscado vivienda en renta en Miahuatlán de Porfirio Díaz.
    \item \textbf{Variables:} el medio de búsqueda (independiente); el tiempo de \nuevo{localización de ofertas} y la confianza percibida \nuevo{en el medio} (dependientes)\nuevo{; la exactitud de la elección y la confianza en los arrendadores (dependientes secundarias)}.
    \item \textbf{Relación lógica:} el medio de búsqueda \emph{reduce} el tiempo de \nuevo{localización de ofertas} y \emph{aumenta} la confianza percibida \nuevo{en el medio}.
\end{itemize}

\subsection{Variables de investigación}

\nuevo{En la Tabla~\ref{tab:operacionalizacion} se define cada variable y la forma en que se medirá o controlará.}

\begin{table}[H]
\centering
\footnotesize
\renewcommand{\baselinestretch}{1}\selectfont % interlineado sencillo en la tabla
\caption[Operacionalización de las variables.]{Operacionalización de las variables.}
\label{tab:operacionalizacion}
\begin{tabularx}{\textwidth}{>{\RaggedRight\arraybackslash}p{2.3cm} >{\RaggedRight\arraybackslash}X >{\RaggedRight\arraybackslash}X >{\RaggedRight\arraybackslash}p{2.7cm}}
\toprule
\textbf{Variable} & \textbf{Definición conceptual} & \textbf{Definición operacional} & \textbf{Indicador e instrumento} \\
\midrule
\multicolumn{4}{l}{\nuevo{\textit{Independiente}}} \\
\nuevo{Medio de búsqueda} &
\nuevo{Canal mediante el cual el estudiante consulta la oferta de vivienda en renta.} &
\nuevo{(a) La plataforma propuesta, con ofertas verificadas; (b) los medios actuales, replicados en un grupo privado de Facebook con las publicaciones originales y fotografías de los anuncios físicos.} &
\nuevo{Condición asignada en cada tarea (nominal).} \\
\addlinespace
\multicolumn{4}{l}{\nuevo{\textit{Dependientes}}} \\
Tiempo de \nuevo{localización de ofertas} &
Lapso que una persona invierte desde que comienza a \nuevo{consultar un medio} hasta que identifica una oferta que satisface sus criterios. \nuevo{No incluye contactar al arrendador, visitar el inmueble ni cerrar el trato.} &
Tiempo que tarda cada participante en completar una tarea estandarizada (localizar una vivienda que cumpla un presupuesto, servicios y distancia máxima a la UNSIS), desde el inicio de la tarea hasta que señala la oferta elegida. \nuevo{Límite de 10 minutos por tarea.} &
Segundos por tarea. Bitácora con cronómetro y grabación de pantalla. \\
\addlinespace
Confianza percibida \nuevo{en el medio} &
\nuevo{Creencia de que el medio cuenta con mecanismos (verificación, información estructurada y actualizada) que hacen fiable lo que publica: confianza basada en la institución \parencite{McKnight2002, Pavlou2004}.} &
\nuevo{Escala adaptada de la confianza basada en la institución \parencite{McKnight2002} y de la efectividad percibida de los mecanismos \parencite{Pavlou2004}, referida al medio. Likert de 1 a 7.} &
Puntuación media de 1 a 7. Cuestionario \nuevo{anónimo} aplicado al terminar la tarea en cada medio. \\
\addlinespace
\nuevo{Confianza en los arrendadores (secundaria)} &
\nuevo{Creencia en la honestidad de los arrendadores que publican en el medio.} &
\nuevo{Subescala breve adaptada de \textcite{Pavlou2004}, en el mismo cuestionario.} &
\nuevo{Puntuación media de 1 a 7.} \\
\addlinespace
\nuevo{Exactitud de la elección (secundaria)} &
\nuevo{Grado en que la oferta elegida satisface los criterios de la tarea.} &
\nuevo{Si la oferta señalada cumple todos los criterios (acierto) y si la tarea se concluyó en el tiempo límite.} &
\nuevo{Proporción de aciertos y tasa de tareas concluidas.} \\
\addlinespace
\multicolumn{4}{l}{\nuevo{\textit{Control}}} \\
\nuevo{Conjunto de ofertas y tareas} &
\nuevo{Material de búsqueda equivalente en ambas condiciones.} &
\nuevo{Dos subconjuntos de ofertas y dos tareas equivalentes, cruzados con el medio y el orden.} &
\nuevo{Diseño del material.} \\
\addlinespace
\nuevo{Entorno de prueba} &
\nuevo{Condiciones físicas de la sesión.} &
\nuevo{Mismo equipo de cómputo, conexión, guion y moderador.} &
\nuevo{Protocolo de sesión.} \\
\addlinespace
\nuevo{Experiencia previa} &
\nuevo{Familiaridad del participante con la búsqueda de vivienda y con Facebook.} &
\nuevo{Número de búsquedas de vivienda realizadas y frecuencia de uso de Facebook.} &
\nuevo{Cuestionario inicial; covariables en el análisis.} \\
\bottomrule
\end{tabularx}
\end{table}

\nuevo{La confianza se mide con las medidas de \textcite{McKnight2002}, construidas para el comercio electrónico, porque distinguen la confianza en el medio de la confianza en la contraparte. No se emplean escalas de confianza en sistemas automatizados, como la de \textcite{Jian2000}, porque provienen del estudio de la automatización \parencite{LeeSee2004} y sus reactivos se refieren a la confiabilidad de un sistema, y un grupo de Facebook no es un sistema automatizado. En este proyecto, la variable de la hipótesis es la confianza en el medio, es decir, en que el canal ofrece garantías sobre la veracidad de lo que publica, porque es lo que la plataforma modifica directamente mediante la verificación; según \textcite{Pavlou2004}, esa confianza institucional se transfiere a la comunidad de oferentes. Para comprobar esta transferencia sin alargar el cuestionario, la confianza en los arrendadores se mide como variable secundaria con una subescala breve y no forma parte del contraste de $H_1$. Los reactivos se traducirán al español y se adaptarán al contexto de la renta; su validez de contenido se evaluará con un panel de expertos mediante la razón de validez de contenido \parencite{Lawshe1975}, y su consistencia interna (alfa de Cronbach) se estimará en una prueba piloto con estudiantes que no formen parte de la muestra.}

\subsection{Hipótesis estadísticas}

Sean $\mu_{T,P}$ y $\mu_{T,M}$ los tiempos medios de \nuevo{localización de ofertas} obtenidos por los participantes con la plataforma y con los medios actuales, respectivamente, y $\mu_{C,P}$ y $\mu_{C,M}$ las puntuaciones medias de confianza \nuevo{en el medio} con la plataforma y con los medios actuales. Puesto que la hipótesis de investigación afirma que la plataforma produce ambos efectos de manera simultánea, la hipótesis alternativa se expresa como la conjunción de las dos condiciones y la hipótesis nula como su negación:

\begin{itemize}
    \item $H_1$: la plataforma reduce el tiempo de \nuevo{localización de ofertas} \emph{y} aumenta la confianza percibida \nuevo{en el medio} frente a los medios actuales
    \[
        H_1:\; (\mu_{T,P} < \mu_{T,M}) \;\wedge\; (\mu_{C,P} > \mu_{C,M}).
    \]
    \item $H_0$: la plataforma no reduce el tiempo de \nuevo{localización de ofertas} \emph{o} no aumenta la confianza percibida \nuevo{en el medio} frente a los medios actuales
    \[
        H_0:\; (\mu_{T,P} \geq \mu_{T,M}) \;\vee\; (\mu_{C,P} \leq \mu_{C,M}).
    \]
\end{itemize}

\subsection[Resumen del diseño de la evaluación]{\nuevo{Resumen del diseño de la evaluación}}

\nuevo{El detalle del procedimiento se desarrollará en el capítulo de metodología; aquí se resumen las decisiones que sostienen la validez del contraste.}

\nuevo{\textbf{Material equivalente y sin arrastre.} Si las mismas viviendas y la misma tarea se presentaran en ambos medios, el participante recordaría en el segundo la respuesta encontrada en el primero; el contrabalanceo equilibra ese sesgo, pero no lo elimina. Por ello, se construirán dos subconjuntos de ofertas (A y B) y dos tareas equivalentes, y se cruzarán con el medio y con el orden, de modo que cada participante busque en un subconjunto distinto en cada medio. Las cuatro combinaciones de orden y subconjunto se asignarán al azar.}

\nuevo{\textbf{Condición de medios actuales.} Para que la comparación sea fiel, las publicaciones originales de Facebook se replicarán, con autorización de los arrendadores, en un grupo privado de Facebook en el que el participante podrá desplazarse y usar el buscador como lo haría normalmente; los anuncios físicos se presentarán como fotografías dentro del mismo grupo.}

\nuevo{\textbf{Mecanismo de verificación.} Se propone un mecanismo de verificación en dos niveles, que se ajustará con los resultados del objetivo específico 1 y se especificará en el objetivo 2: (a) \emph{verificación del arrendador}, en la que las tesistas cotejan en persona su identificación oficial y confirman su teléfono, sin almacenar copia del documento; y (b) \emph{verificación del inmueble}, mediante una visita en la que se confirma su existencia, se toman fotografías georreferenciadas y se comprueba que las características publicadas (precio, servicios y reglas) coinciden con las reales. La plataforma mostrará una insignia distinta para cada nivel, la fecha de la última verificación y un botón para que los usuarios reporten anuncios incorrectos. Todas las ofertas que aparezcan como verificadas en la evaluación habrán pasado efectivamente por este mecanismo, para que la confianza medida responda al mecanismo y no solo a una etiqueta.}

\nuevo{\textbf{Control de la deseabilidad social.} Los cuestionarios serán anónimos, se seguirá un guion estandarizado y, de ser posible, las sesiones las conducirá un moderador ajeno al desarrollo de la plataforma.}

\nuevo{\textbf{Tiempo límite.} Las tareas no concluidas en 10 minutos se registrarán con el tiempo límite. Se reportará cuántos casos lo alcanzan y se realizará un análisis de sensibilidad excluyéndolos, para valorar el efecto de techo.}
% TODO: confirmar el tiempo límite tras la prueba piloto.

\nuevo{\textbf{Tamaño de muestra.} Para una prueba $t$ pareada unilateral con $\alpha = 0.05$, potencia de 0.80 y un efecto mediano ($d = 0.5$; \cite{Cohen1988}), se requieren 27 participantes \parencite{Faul2007}. Como la prueba de intersección-unión exige que ambas pruebas resulten significativas, se busca una potencia de 0.90 en cada una, lo que eleva el requerimiento a 36 participantes; considerando un 10\,\% de pérdidas, se reclutarán al menos 40 estudiantes no residentes.}
% Cálculo verificado con la distribución t no central (n = 27 con potencia 0.81; n = 36 con potencia 0.90). Ajustar d si la prueba piloto sugiere otro valor.

\subsection[Análisis estadístico]{\nuevo{Análisis estadístico}}

\nuevo{Cada} una de las dos condiciones \nuevo{de $H_1$} se evaluará con una prueba $t$ de Student para muestras relacionadas, unilateral y con $\alpha = 0.05$. \nuevo{Dado que los tiempos suelen tener una distribución asimétrica, el tiempo de localización se analizará tras una transformación logarítmica, definida de antemano; la prueba de rangos con signo de Wilcoxon se reportará como análisis de robustez.} La decisión se tomará mediante una prueba de intersección-unión: $H_0$ se rechazará únicamente si las dos pruebas unilaterales resultan significativas, lo que conserva el nivel de significancia global en $\alpha = 0.05$ sin requerir correcciones por comparaciones múltiples \parencite{BergerHsu1996}. Si solo una de las condiciones resulta significativa, la hipótesis de investigación no se considerará confirmada, aunque el efecto observado sobre esa variable se reportará y discutirá.

\nuevo{Además del valor $p$, se reportarán la diferencia de medias con su intervalo de confianza al 95\,\% y la $d$ de Cohen. La experiencia previa en la búsqueda de vivienda y la frecuencia de uso de Facebook se incorporarán como covariables en un modelo lineal mixto con el medio, el orden y el subconjunto de ofertas como efectos fijos y el participante como efecto aleatorio. La exactitud de la elección y la tasa de tareas concluidas se compararán entre medios con la prueba de McNemar y se reportarán de forma descriptiva. La confianza en los arrendadores se comparará con una prueba $t$ pareada como análisis secundario, sin intervenir en la decisión sobre $H_0$.}

\subsection[Limitaciones del diseño]{\nuevo{Limitaciones del diseño}}

\begin{itemize}
    \item \nuevo{\textbf{Efecto de paquete.} La variable independiente combina centralización, geolocalización y verificación; el diseño permite saber si la plataforma en conjunto produce el efecto, pero no cuál de sus componentes lo causa.}
    \item \nuevo{\textbf{Alcance del tiempo medido.} Se mide el tiempo para localizar una oferta en pantalla, no el tiempo total para conseguir vivienda, que incluye contactar al arrendador, visitar el inmueble y cerrar el trato.}
    \item \nuevo{\textbf{Situación simulada.} Las tareas se realizan en una sesión controlada y no en una búsqueda real, lo que limita la generalización de los resultados.}
\end{itemize}

\subsection{Medición complementaria: usabilidad}

La usabilidad no forma parte de la hipótesis, puesto que no es una de las variables cuya relación se contrasta, pero se medirá como evaluación complementaria en atención al \nuevo{sexto} objetivo específico. \nuevo{Los estudiantes responderán la Escala de Usabilidad del Sistema \parencite[SUS;][]{Brooke1996} al concluir las tareas; los arrendadores la responderán tras realizar tareas propias de su rol (registrar un inmueble, completar la verificación y actualizar su disponibilidad). La escala consta de 10 reactivos y} produce una puntuación de 0 a 100. El resultado se reportará de forma descriptiva y se interpretará con respecto al promedio de referencia de 68 puntos \parencite{LewisSauro2018} y a la escala de adjetivos de \textcite{Bangor2008}. Este dato también ayudará a interpretar los resultados del contraste de la hipótesis; por ejemplo, una baja usabilidad podría explicar que la plataforma no reduzca el tiempo de \nuevo{localización}.

\section[Matriz de congruencia]{\nuevo{Matriz de congruencia}}

\nuevo{La Tabla~\ref{tab:congruencia} relaciona cada pregunta específica con su objetivo, su criterio de logro y la técnica con la que se obtendrá la evidencia.}

% TODO: sustituir los corchetes por los tamaños definitivos.
\begin{table}[H]
\centering
\footnotesize
\renewcommand{\baselinestretch}{1}\selectfont % interlineado sencillo en la tabla
\caption[Matriz de congruencia.]{\nuevo{Matriz de congruencia.}}
\label{tab:congruencia}
\begin{tabularx}{\textwidth}{>{\RaggedRight\arraybackslash}X >{\RaggedRight\arraybackslash}X >{\RaggedRight\arraybackslash}X >{\RaggedRight\arraybackslash}X}
\toprule
\textbf{\nuevo{Pregunta específica}} & \leavevmode\textbf{\nuevo{Objetivo específico}} & \leavevmode\textbf{\nuevo{Hipótesis o criterio de logro}} & \leavevmode\textbf{\nuevo{Técnica y evidencia}} \\
\midrule
\nuevo{¿Qué necesidades, criterios y riesgos reportan estudiantes y arrendadores?} & \nuevo{1. Identificar necesidades, criterios de búsqueda y riesgos} & \nuevo{Lista priorizada de requisitos} & \nuevo{Encuesta a [\textit{n}] estudiantes; entrevistas a [\textit{n}] arrendadores} \\
\addlinespace
\nuevo{¿Qué mecanismos de verificación son viables?} & \nuevo{2. Diseñar arquitectura, modelo de datos y verificación} & \nuevo{Especificación del mecanismo aceptada por ambas partes} & \nuevo{Revisión de literatura; prototipos; validación con usuarios} \\
\addlinespace
\nuevo{¿El prototipo cumple los requisitos?} & \nuevo{3. Implementar el prototipo con registro, verificación, búsqueda y mapa} & \nuevo{Requisitos funcionales verificados} & \nuevo{Pruebas funcionales} \\
\addlinespace
\nuevo{¿Con qué ofertas se evalúa?} & \nuevo{4. Construir el conjunto de ofertas reales verificadas} & \nuevo{[\textit{n}] ofertas autorizadas, geolocalizadas y verificadas} & \nuevo{Recopilación en campo; consentimiento de arrendadores} \\
\addlinespace
\nuevo{¿La plataforma reduce el tiempo y aumenta la confianza?} & \nuevo{5. Comparar tiempo, exactitud y confianza entre medios} & \nuevo{$H_1$: menor tiempo y mayor confianza con la plataforma} & \nuevo{Experimento intrasujetos con al menos 40 participantes; $t$ pareada (Wilcoxon como robustez); intersección-unión, $\alpha = 0.05$} \\
\addlinespace
\nuevo{¿Qué usabilidad percibida alcanza?} & \nuevo{6. Evaluar la usabilidad con estudiantes y arrendadores} & \nuevo{SUS igual o mayor que 68} & \nuevo{Escala SUS; estadística descriptiva} \\
\bottomrule
\end{tabularx}
\end{table}




    %--- Apéndices ---
    \appendix
 %   \chapter{Título del Apéndice}
% Escribe aquí el contenido del apéndice...
  % Apéndice A
    % \include{capitulos/apendiceB}  % Apéndice B 

    \printbibliography[title={Referencias}]

\end{document}
